% ==============================================================================
% Front matter: Group Work Allocation (VNU-UET style)
% ==============================================================================

\chapter*{Phân công công việc nhóm}
\phantomsection
\addcontentsline{toc}{chapter}{Phân công công việc nhóm}

\draftnote[Chung]{Cập nhật lại họ tên, MSSV và tỷ lệ đóng góp thực tế của từng thành viên trước khi nộp báo cáo.}

\begin{table}[H]
\centering
\small
\renewcommand{\arraystretch}{1.25}
\begin{tabularx}{\textwidth}{@{}P{2.6cm}P{3.4cm}L M{2.2cm}@{}}
\toprule
\rowcolor{tblhead}
\thead{Thành viên} & \thead{Bài báo làm chủ} & \thead{Nội dung báo cáo phụ trách} & \thead{Ước lượng} \\
\midrule
\textbf{Thành viên 1}\newline {\footnotesize(Điều phối)} &
\textit{Ask in Any Modality} (ACL '25)\newline \textit{UniversalRAG} (ACL '26) &
Chương 1; Chương 2; Mục 3.1 (\textit{UniversalRAG}); phần dữ liệu, độ đo, bảng số liệu và phản biện \textit{UniversalRAG} trong Chương 4; ghép bài và kiểm tra định dạng chung. &
14--16 trang \\
\midrule
\textbf{Thành viên 2} &
\textit{ColPali} (ICLR '25)\newline \textit{VisRAG} (ICLR '25) &
Toàn bộ Mục 3.2 (\textit{ColPali}, \textit{VisRAG} và đối chứng Mục 3.2.4); phần benchmark ViDoRe/VQA/BEIR, so sánh trực tiếp ColPali--VisRAG và phản biện 2 bài trong Chương 4; chuẩn hóa sơ đồ Chương 3. &
10--12 trang \\
\midrule
\textbf{Thành viên 3} &
\textit{CoRe-MMRAG} (ACL '25)\newline Tích hợp hệ thống &
Mục 3.3 (\textit{CoRe-MMRAG}); Mục 3.4 (so sánh 4 bài \& kiến trúc tích hợp); phần InfoSeek/Enc-VQA, Mục 4.3.1, bảng số liệu và phản biện \textit{CoRe-MMRAG} trong Chương 4; Chương 6; rà soát \texttt{references.bib}. &
11--13 trang \\
\bottomrule
\end{tabularx}
\end{table}
